% Part: lambda-calculus
% Chapter: lambda-definability
% Section: minimization

\documentclass[../../../include/open-logic-section]{subfiles}

\begin{document}

\olfileid{lam}{ldf}{min}
\olsection{लघुतमीकरण}

मूलभूत फलने $\Zero$, $\Succ$, $\Proj{n}{i}$
यांपासून संयोजन, आदिम पुनरावर्तन आणि नियमित
लघुतमीकरण वापरून मिळणारी फलने म्हणजे सामान्य
पुनरावर्ती फलने. सर्व सामान्य पुनरावर्ती फलने
!!{lambda definable} आहेत हे दाखवण्यासाठी,
!!a{lambda
  definable} फलनापासून नियमित लघुतमीकरणाने
परिभाषित केलेले कोणतेही फलनही !!{lambda definable}
आहे हे दाखवावे लागेल.

\begin{lem}
  \ollabel{lem:min} $f(x_1, \dots, x_k, y)$ हे नियमित
  आणि !!{lambda definable} असेल, तर
  \[
  h(x_1, \dots, x_k) = \umin{y}{f(x_1,\dots,x_k, y) = 0}
  \]
  याने परिभाषित केलेले $h$~देखील
  !!{lambda definable} आहे.
\end{lem}

\begin{proof}
  लॅम्डा पद~$F$ हे नियमित फलन $f(\vec x, y)$ ला
  $\lambda$-परिभाषित करते असे समजा. $h$ ला
  !!{lambda define} करण्यासाठी शोधफलन आणि
  स्थिरबिंदू संयोजक वापरू:
  \begin{align*}
    \fn{Search} & \ident \lambd[g][\lambd[f\,\vec{x}\,y][
        \fn{IsZero} (f\, \vec{x}\, y)\, y\, (g\, \vec{x} (\fn{Succ}\, y))]]\\
    H & \ident \lambd[\vec x][(Y \, \fn{Search}) F\, \vec{x}\, \num{0}],
  \end{align*}
  येथे $Y$ हा कोणताही स्थिरबिंदू संयोजक आहे.
  सहज अर्थाने, $\fn{Search}$ हे स्वतःला पुन्हा
  बोलावणारे फलन आहे: $y$ पासून सुरू करून
  $f\, \vec x\, y$ शून्य आहे का हे तपासते;
  शून्य असेल तर~$y$ परत करते, अन्यथा
  $\fn{Succ}\, y$ घेऊन स्वतःला पुन्हा बोलावते.
  म्हणून $(Y \, \fn{Search}) F
  \num{n_1}\dots\num{n_k}\,\num{0}$ हे
  $f(n_1, \dots, n_k, m) = 0$ पूर्ण करणारा
  लघुतम~$m$ परत करते.

  नेमके सांगायचे तर, पुढील गोष्ट पाहा:
  \begin{align*}
    (Y \, \fn{Search}) F \num{n_1}
    \dots\num{n_k}\, \num{m} & \red \num{m}
    \intertext{जर $f(n_1, \dots,
      n_k, m) = 0$ असेल; अन्यथा}
    & \red (Y \, \fn{Search}) F\, \num{n_1} \dots\num{n_k}\, \num{m+1}
    \intertext{मिळते. $f$ नियमित असल्याने काही $y$ साठी
      $f(n_1, \dots, n_k, y) = 0$ असते; म्हणून}
    (Y \, \fn{Search}) F \num{n_1} \dots\num{n_k}\,\num{0}
    & \red \num{h(n_1, \dots, n_k)}.
    \end{align*}
\end{proof}

\begin{prop}
  प्रत्येक सामान्य पुनरावर्ती फलन !!{lambda definable} आहे.
\end{prop}

\begin{proof}
 \olref[prf]{lem:basic} नुसार सर्व मूलभूत फलने
 !!{lambda definable} आहेत. तसेच \olref[prf]{lem:comp},
 \olref[prf]{lem:prim} आणि \olref{lem:min} नुसार
 !!{lambda definable} फलनांचा वर्ग संयोजन,
 आदिम पुनरावर्तन आणि नियमित लघुतमीकरण यांच्या
 अंतर्गत बंद असतो.
\end{proof}

\end{document}
